Ce travail porte sur la saisie autonome d'objets par un bras robotique à bas coût guidé par la vision. Sur un SO-101 équipé de trois caméras --- une paire stéréo \emph{eye-to-hand} et une caméra \emph{eye-in-hand} montée au poignet ---, nous construisons une \emph{pipeline} complète et modulaire allant du pixel à la prise : détection de l'objet, reconstruction de sa position tridimensionnelle, planification de la saisie, puis contrôle du bras. L'ensemble est \emph{open-source} et reproductible, et chacune de ses étapes est conçue et expliquée séparément plutôt que traitée comme une boîte noire.

La planification de saisie règle l'ouverture et l'orientation des mâchoires sur la géométrie mesurée de l'objet et choisit un angle d'attaque par balayage. Pour mettre cette approche explicite en perspective, nous implémentons sur le \emph{même} bras et la \emph{même} tâche une seconde méthode --- une politique apprise de bout en bout par imitation (\textit{Action Chunking Transformer}) ---, puis nous comparons les deux paradigmes sur un banc commun.

Deux campagnes de mesures --- 78 essais pour la pipeline, 30 pour la politique apprise, assorties d'intervalles de confiance de Wilson --- établissent des taux de réussite de 67 à 77\,\% pour la pipeline sur le cube et de 67 à 78\,\% sur le cylindre (selon le détecteur employé), et de 70\,\% pour la politique apprise sur le cube. À égalité sur l'objet entraîné, la pipeline l'emporte en généralité et en transparence, la politique apprise en rapidité de mise en œuvre. Le travail assume enfin ses limites : un bras à cinq degrés de liberté --- soit un de moins que les six d'une pose quelconque dans l'espace ---, l'absence délibérée d'évitement d'obstacles, et un périmètre d'objets et de conditions explicitement déclaré.
